% ===============================================================
%  جلسه چهارم
% ===============================================================
\chapter{احتمال شرطی، قانون زنجیره‌ای، احتمال کل و قانون بیز}

\section{احتمال شرطی \lr{(Conditional Probability)}}

در بسیاری موارد، \textbf{اطلاعاتی} در مورد یک واقعه تصادفی به دست ما می‌رسد. در این شرایط، به دنبال یافتن احتمال‌های به‌روزشده هستیم.

برای مثال، در پرتاب یک تاس ۶ وجهی متقارن، در صورتی که پیش‌تر بدانیم واقعه \textbf{آمدن یک عدد فرد}، یعنی $E_{\text{odd}}=\{1,3,5\}$، حتماً اتفاق افتاده است، آنگاه:
\begin{itemize}
\item احتمال رخ دادن عدد ۲ \textbf{صفر} خواهد شد (پیش از آن $\frac16$ بود)؛
\item احتمال واقعه رخ دادن عددی کوچک‌تر یا مساوی ۳ \textbf{بیشتر} خواهد شد: پیش از اطلاع برابر $\frac36=\frac12$ بود، اما پس از دانستن فرد بودن، از سه حالت ممکن $\{1,3,5\}$ دو حالت $\{1,3\}$ مطلوب‌اند؛ پس احتمال به $\frac23$ می‌رسد.
\end{itemize}
این در حالی است که این احتمالات پیش از رسیدن اطلاعات فوق، اعداد متفاوتی بودند.

\begin{definition}[احتمال شرطی]
در آمار و احتمال، \textbf{احتمال شرطی} واقعه $A$ به شرط $B$ (خوانده می‌شود \lr{$A$ given $B$} یا \lr{$A$ conditioned on $B$}) را به صورت احتمال به‌روزشده رویداد $A$ تعریف می‌کنیم، به شرطی که بدانیم رویداد $B$ حتماً اتفاق افتاده است. برای $\Prob(B)>0$ تعریف می‌کنیم:
\[
\Prob(A\given B)=\frac{\Prob(A\cap B)}{\Prob(B)} .
\]
\end{definition}

\textbf{دلیل ظاهر شدن مخرج:} پس از دانستن رخ دادن $B$، گویی فضای نمونه‌ها $\Omega$ کوچک‌تر شده است. تنها وقایعی که برابر یا زیرمجموعه $B$ هستند باقی خواهند ماند. پس از $A$ فقط قسمت $A\cap B$ هنوز ممکن است، و برای اینکه احتمال کل فضای جدید (یعنی $B$) برابر ۱ شود، همه احتمال‌ها را بر $\Prob(B)$ تقسیم می‌کنیم (نرمال‌سازی). در مثال تاس:
\[
\Prob(\le3\given \text{odd})=\frac{\Prob(\{1,3\})}{\Prob(\{1,3,5\})}=\frac{2/6}{3/6}=\frac23 .
\]

\begin{figure}[H]
\centering
\begin{tikzpicture}[scale=1.1]
\fill[red!25] (0,0) ellipse (3 and 1.6);
\draw[thick] (0,0) ellipse (3 and 1.6);
\fill[white] (-0.4,0.1) circle (0.75);
\fill[gray!45] (0.4,0.1) circle (0.75);
\draw[thick] (-0.4,0.1) circle (0.75);
\draw[thick] (0.4,0.1) circle (0.75);
\node[green!50!black] at (-0.75,0.1) {\Large $A$};
\node[green!50!black] at (0.8,0.1) {\Large $B$};
\node at (-2.6,1.4) {$\Omega$};
\begin{scope}[xshift=4.1cm]
\fill[red!25] (0,0.7) rectangle (0.4,1);
\node[anchor=west,font=\small,align=left] at (0.5,0.85) {ناممکن (زیرا $B$ رخ داده است)};
\fill[gray!45] (0,-0.3) rectangle (0.4,0);
\node[anchor=west,font=\small,align=left] at (0.5,-0.15) {فضای نمونه جدید پس از رخ دادن $B$};
\end{scope}
\end{tikzpicture}
\caption{نمایش گرافیکی احتمال شرطی: بعد از رخداد $B$، فضای نمونه مؤثر (و به تبع آن، فضای وقایع مؤثر) کوچک‌تر خواهد شد.}
\end{figure}

\begin{note}
\notetitle دقت کنید که $\Prob(A\given B)$ می‌تواند از $\Prob(A)$ \textcolor{red!70!black}{\textbf{بزرگ‌تر}}، \textcolor{blue!70!black}{\textbf{کوچک‌تر}} یا \textcolor{green!50!black}{\textbf{مساوی}} باشد. این بستگی خواهد داشت که رویداد $B$، چه میزان اطلاعاتی در مورد رویداد $A$ به دست می‌دهد. در مثال تاس: $\Prob(\le3\given\text{odd})=\frac23>\frac12$ ولی $\Prob(2\given\text{odd})=0<\frac16$.
\end{note}

\begin{theorem}[احتمال شرطی و استقلال]
اگر $\Prob(B)>0$، آنگاه دو رویداد $A$ و $B$ مستقل‌اند اگر و تنها اگر $\Prob(A\given B)=\Prob(A)$؛ زیرا در صورت استقلال
\[
\Prob(A\given B)=\frac{\Prob(A\cap B)}{\Prob(B)}=\frac{\Prob(A)\Prob(B)}{\Prob(B)}=\Prob(A),
\]
و برعکس، از $\Prob(A\given B)=\Prob(A)$ با ضرب دو طرف در $\Prob(B)$ تعریف استقلال به دست می‌آید.
\end{theorem}
این تعبیر، استقلال را شهودی‌تر می‌کند: \emph{دانستن رخ دادن $B$ هیچ اطلاعی درباره $A$ نمی‌دهد}.

\begin{note}
\notetitle به ازای $B$ ثابت، تابع $A\mapsto\Prob(A\given B)$ خودش یک اندازه احتمال (در اصول کولموگروف صادق) است: نامنفی است، $\Prob(\Omega\given B)=\Prob(B)/\Prob(B)=1$، و جمع‌پذیر است. پس همه قضایایی که برای احتمال ثابت کردیم (متمم، شمول و عدم شمول، کران اجتماع و \ldots) برای احتمال شرطی نیز برقرارند؛ مثلاً $\Prob(A^c\given B)=1-\Prob(A\given B)$.
\end{note}

\begin{example}[تصادفات جاده‌ای و شرایط جوی]
فرض کنید احتمالات زیر مربوط به واقعه تصادفات جاده‌ای در شرایط مختلف جوی باشند:
\begin{center}
\renewcommand{\arraystretch}{1.4}
\begin{tabular}{c|c|c}
 & صاف (\lr{clear}) & مه‌آلود (\lr{mist})\\\hline
تردد امن & $0.73$ & $0.17$\\\hline
تصادف (\lr{crash}) & $0.01$ & $0.09$
\end{tabular}
\end{center}
\begin{itemize}
\item سؤال ۱: احتمال تصادف چقدر است؟
\item سؤال ۲: احتمال تصادف به شرط هوای مه‌آلود چقدر است؟
\item سؤال ۳: آیا احتمال تصادف، از شرایط جوی جاده مستقل است؟
\end{itemize}
\begin{solution}
\textbf{سؤال ۱:} رویداد تصادف اجتماع دو خانه ناسازگار سطر دوم است:
\[\Prob(\text{crash})=0.01+0.09=0.10 .\]
\textbf{سؤال ۲:} ابتدا احتمال مه‌آلود بودن (جمع ستون مه‌آلود) و احتمال اشتراک را می‌یابیم:
\[
\Prob(\text{crash}\cap\text{mist})=0.09,\qquad \Prob(\text{mist})=0.17+0.09=0.26,
\]
\[
\Prob(\text{crash}\given\text{mist})=\frac{0.09}{0.09+0.17}=\frac{0.09}{0.26}\approx0.346 .
\]
\textbf{سؤال ۳:} $\Prob(\text{crash}\given\text{mist})\approx0.346$ با $\Prob(\text{crash})=0.10$ \textbf{برابر نیستند}. پس واقعه تصادف مستقل از مه‌آلود بودن نیست. به‌علاوه در هوای صاف $\Prob(\text{crash}\given\text{clear})=\frac{0.01}{0.74}\approx0.0135$؛ یعنی احتمال تصادف در هوای مه‌آلود حدود ۲۶ برابر هوای صاف است!
\end{solution}
\end{example}

\section{قانون زنجیره‌ای \lr{(Chain Rule)}}

\subsection{حالت دوتایی و سه‌تایی}
با توجه به تعریف احتمال شرطی (و نوشتن $\Prob(A,B)$ به جای $\Prob(A\cap B)$) خواهیم داشت:
\[
\boxed{\ \Prob(A,B)=\Prob(A\given B)\,\Prob(B)\ }
\]
در صورتی که سه واقعه $A$، $B$ و $C$ را مدنظر قرار داده باشیم، می‌توان رابطه بالا را به صورت تو در تو تکرار کرد. ابتدا $(B,C)$ را یک رویداد در نظر می‌گیریم:
\[
\Prob(A,B,C)=\Prob(A\given B,C)\,\Prob(B,C),
\]
و سپس رابطه دوتایی را برای $\Prob(B,C)$ به کار می‌بریم:
\[
\boxed{\ \Prob(A,B,C)=\Prob(A\given B,C)\,\Prob(B\given C)\,\Prob(C)\ }
\]

\begin{note}
\notetitle وقتی احتمالی را به بیش از یک رویداد مشروط می‌کنیم (مانند $\Prob(A\given B,C)$)، منظور این است که فرض کنید \emph{تمامی} آنان رخ داده باشند؛ یعنی $\Prob(A\given B,C)=\Prob(A\given B\cap C)$.
\end{note}

\subsection{حالت کلی و موارد کاربردی}

\begin{theorem}[قانون زنجیره‌ای]
در حالت کلی، برای احتمال رخ دادن توأمان تمامی رویدادهای $A_1,A_2,\ldots,A_n$ خواهیم داشت:
\[
\Prob(A_1,A_2,\ldots,A_n)=\Prob(A_1)\prod_{i=2}^{n}\Prob(A_i\given A_{i-1},\ldots,A_1).
\]
\end{theorem}
\begin{proof}
با استقرا روی $n$. برای $n=2$ همان تعریف احتمال شرطی است. اگر برای $n-1$ برقرار باشد، با در نظر گرفتن $B=A_1\cap\cdots\cap A_{n-1}$ داریم
$\Prob(A_1,\ldots,A_n)=\Prob(A_n\given B)\Prob(B)$ و با جایگذاری فرض استقرا برای $\Prob(B)$ حکم ثابت می‌شود. (معادلاً، در حاصل‌ضرب سمت راست، صورت هر کسر با مخرج کسر بعدی ساده می‌شود و فقط $\Prob(A_1\cap\cdots\cap A_n)$ باقی می‌ماند.)
\end{proof}

در مواردی محاسبه احتمالات شرطی بسیار ساده است. مثلاً، یک مدل‌سازی آماری را از ابتدا بر حسب احتمال‌های شرطی بنا می‌کنیم (مانند مسئله روز تولد در جلسه قبل که در آن $\Prob(E_i\given E_2,\ldots,E_{i-1})=1-\frac{i-1}{365}$ به‌راحتی به دست می‌آمد).

\begin{note}
\notetitle \textbf{نکته اضافی (خاصیت مارکوف):} بسیاری از پدیده‌ها، خاصیت \lr{Markovian} دارند. به عنوان مثال، رویداد $A_i$ می‌تواند به شرط دانستن $A_{i-1}$ از تمام رویدادهای $A_{i-2},A_{i-3},\ldots,A_1$ مستقل شود. در این صورت، محاسبه رابطه فوق بسیار ساده می‌شود:
\[
\Prob(A_1,\ldots,A_n)=\Prob(A_1)\prod_{i=2}^{n}\Prob(A_i\given A_{i-1}).
\]
برای مثال آب‌وهوای فردا را معمولاً تنها با دانستن آب‌وهوای امروز پیش‌بینی می‌کنیم.
\end{note}

\begin{example}[بیرون کشیدن بدون جایگذاری]
از یک دسته کارت ۵۲تایی، سه کارت بدون جایگذاری بیرون می‌کشیم. احتمال اینکه هر سه کارت آس باشند چقدر است؟
\begin{solution}
$A_i$ را رویداد «کارت $i$ام آس است» بگیرید. با قانون زنجیره‌ای:
\[
\Prob(A_1,A_2,A_3)=\Prob(A_1)\Prob(A_2\given A_1)\Prob(A_3\given A_1,A_2)=\frac4{52}\cdot\frac3{51}\cdot\frac2{50}=\frac{24}{132600}\approx1.81\times10^{-4}.
\]
هر احتمال شرطی به‌سادگی به دست می‌آید: پس از بیرون آمدن یک آس، ۳ آس در ۵۱ کارت باقی مانده است. (بررسی با شمارش: $\binom43/\binom{52}3=4/22100$، همان عدد.)
\end{solution}
\end{example}

\section{قانون احتمال کل \lr{(The Law of Total Probability)}}

\begin{definition}[افراز]
فرض کنید که $A_1,A_2,\ldots,A_n$ (که $n$ می‌تواند $\infty$ هم باشد) رویدادهایی درون $\F$ باشند. آن‌ها را یک \textbf{افراز} (\lr{Partition}) از $\Omega$ گوییم اگر
\[
\bigcup_{i=1}^{n}A_i=\Omega,\qquad A_i\cap A_j=\varnothing\quad(\text{if } i\neq j).
\]
در این صورت، با هر بار انجام آزمایش یکی (و دقیقاً یکی) از $A_1,\ldots,A_n$ رخ خواهد داد. همچنین تمام $\Prob(A_i)$ها نیز تعریف شده هستند.
\end{definition}

\begin{theorem}[قانون احتمال کل]
اگر $A_1,\ldots,A_n$ یک افراز از $\Omega$ باشند، برای هر واقعه‌ای مانند $E$ خواهیم داشت:
\[
\Prob(E)=\sum_{i=1}^{n}\Prob(E\given A_i)\,\Prob(A_i).
\]
\end{theorem}
\begin{proof}
رویدادهای $E\cap A_i$ ناسازگارند (چون $A_i$ها ناسازگارند) و اجتماعشان $E\cap\Omega=E$ است. پس
$\Prob(E)=\sum_i\Prob(E\cap A_i)=\sum_i\Prob(E\given A_i)\Prob(A_i)$.
\end{proof}

\begin{figure}[H]
\centering
\begin{tikzpicture}[scale=1.0]
\begin{scope}[rotate=35]
\fill[cyan!30] (0,0) ellipse (3 and 1.7);
\fill[olive!55,opacity=0.8] (0,0) ellipse (1.4 and 1.0);
\draw[very thick] (0,0) ellipse (3 and 1.7);
\draw[brown!60!black,thick] (0,0) ellipse (1.4 and 1.0);
\coordinate (O) at (0,0);
\foreach \a in {20,60,110,150,250} { \draw (O) -- ({3*cos(\a)},{1.7*sin(\a)}); }
\end{scope}
\node at (0.95,1.95) {$A_1$};
\node at (1.85,1.0) {$A_2$};
\node at (1.4,-0.1) {$A_3$};
\node at (-0.25,1.75) {$A_n$};
\node at (-1.2,-0.8) {$\cdots$};
\node at (-2.9,1.1) {$\Omega$};
\draw[-{Stealth},thick] (-3.2,2.5) node[above]{\lr{Event} $E$} to[bend right=30] (-0.6,0.2);
\end{tikzpicture}
\caption{قانون احتمال کل: رویداد $E$ (بیضی داخلی) توسط افراز $A_1,\ldots,A_n$ به قطعه‌هایی (\lr{slice}) تقسیم می‌شود.}
\end{figure}

در واقع هر یک از $\Prob(E\given A_i)\Prob(A_i)$ها، طبق تعریف برابر با $\Prob(E,A_i)$ هستند، و برابر با جرم احتمال جمع‌شده در یکی از تکه‌ها (\lr{slice}های) واقعه $E$ می‌باشند.

\begin{note}
\notetitle قانون احتمال کل زمانی به درد می‌خورد که وقوع یک واقعه از طریق چندین \textbf{مکانیزم} وجود دارد، و احتمال وقوع به شرط هر یک از آنان را می‌دانیم. به این روش «شکستن به حالات» (\lr{conditioning}) نیز می‌گویند.
\end{note}

\begin{example}[کارخانه با سه خط تولید]
یک کارخانه سه خط تولید دارد که به ترتیب ۵۰٪، ۳۰٪ و ۲۰٪ محصولات را تولید می‌کنند. درصد محصولات معیوب در این خطوط به ترتیب ۱٪، ۲٪ و ۵٪ است. احتمال اینکه یک محصول تصادفی معیوب باشد چقدر است؟
\begin{solution}
خطوط تولید $L_1,L_2,L_3$ یک افراز هستند. با قانون احتمال کل:
\[
\Prob(D)=0.01\times0.5+0.02\times0.3+0.05\times0.2=0.005+0.006+0.010=0.021 .
\]
\end{solution}
\end{example}

\section{قانون بیز \lr{(Bayes' Rule)}}

قانون بیز، رابطه ساده ولی بسیار پرکاربردی است که کمک می‌کند تا ترتیب وقایع شرطی در احتمال را عوض کنیم. یعنی با داشتن $\Prob(A\given B)$، و همین‌طور احتمال‌های حاشیه‌ای $\Prob(B)$ و $\Prob(A)$، بتوانیم $\Prob(B\given A)$ را محاسبه کنیم.

\begin{theorem}[قانون بیز]
\[
\Prob(B\given A)=\frac{\Prob(B\cap A)}{\Prob(A)}=\frac{\Prob(A\given B)\,\Prob(B)}{\Prob(A)} .
\]
\end{theorem}
\begin{proof}
طبق قانون زنجیره‌ای، $\Prob(A\cap B)$ را به دو صورت می‌توان نوشت:
$\Prob(B\given A)\Prob(A)=\Prob(A\cap B)=\Prob(A\given B)\Prob(B)$. با تقسیم بر $\Prob(A)$ حکم به دست می‌آید.
\end{proof}

یک شکل دیگر نگاه کردن به این رابطه، به صورت زیر است:
\[
\Prob(B\given A)=\frac{\Prob(A\given B)}{\Prob(A)}\,\Prob(B).
\]
در رابطه بالا:
\begin{itemize}
\item به احتمال $\Prob(B)$ در اصطلاح \textbf{احتمال پیشین} (\lr{Prior Probability})، و
\item به احتمال $\Prob(B\given A)$ در اصطلاح \textbf{احتمال پسین} (\lr{Posterior Probability}) می‌گوییم.
\end{itemize}
ضریب $\frac{\Prob(A\given B)}{\Prob(A)}$ نشان می‌دهد مشاهده $A$ باور ما به $B$ را چند برابر می‌کند. به $\Prob(A\given B)$ (به عنوان تابعی از $B$) «درست‌نمایی» (\lr{likelihood}) می‌گویند.

\begin{theorem}[قانون بیز همراه با احتمال کل]
اگر $B_1,\ldots,B_n$ یک افراز باشند، مخرج قانون بیز را می‌توان با قانون احتمال کل حساب کرد:
\[
\Prob(B_j\given A)=\frac{\Prob(A\given B_j)\Prob(B_j)}{\sum_{i=1}^{n}\Prob(A\given B_i)\Prob(B_i)} .
\]
\end{theorem}

\subsection{دیدگاه بیزی}
دیدگاه \lr{Bayesian} در آمار و احتمال بر این اصل مبتنی است که: برای مدل‌سازی تصادفی تأثیر دو رویداد تصادفی بر یکدیگر (مثلاً تأثیر فضای وقایع $\Omega_B$ بر $\Omega_A$)، با فضای وقایع کلی $\Omega_A\times\Omega_B$، صرفاً لازم است تا موارد زیر را داشته باشیم (یا مدل‌سازی کنیم):
\begin{itemize}
\item احتمال در فضای وقایع منتسب به $\Omega_B$ (احتمال پیشین)؛
\item احتمالات شرطی در فضای وقایع منتسب به $\Omega_A$ به شرط وقایع $B$ (مدل مشاهده).
\end{itemize}
با این دو، طبق قانون زنجیره‌ای کل توزیع توأم مشخص می‌شود و با قانون بیز می‌توان «برعکس» فکر کرد: از مشاهده $A$ درباره علت $B$ استنتاج کرد. موارد کاربردی دیدگاه \lr{Bayesian} بسیار زیاد هستند:
\lr{Bayesian inference, Bayesian optimization, Bayesian network, Bayesian Hierarchy} و \ldots

\begin{example}[آزمایش پزشکی ــ یک مثال کلاسیک]
بیماری خاصی در ۱٪ جامعه وجود دارد. یک آزمایش تشخیصی، فرد بیمار را با احتمال ۹۹٪ مثبت اعلام می‌کند (حساسیت) و فرد سالم را با احتمال ۹۵٪ منفی اعلام می‌کند. اگر جواب آزمایش یک فرد تصادفی مثبت باشد، احتمال اینکه واقعاً بیمار باشد چقدر است؟
\begin{solution}
$D$ را رویداد بیماری و $+$ را مثبت شدن آزمایش بگیرید. داریم $\Prob(D)=0.01$، $\Prob(+\given D)=0.99$ و $\Prob(+\given D^c)=0.05$. با احتمال کل:
\[
\Prob(+)=0.99\times0.01+0.05\times0.99=0.0099+0.0495=0.0594 .
\]
با قانون بیز:
\[
\Prob(D\given +)=\frac{0.99\times0.01}{0.0594}\approx0.167 .
\]
یعنی با وجود دقت ظاهراً بالای آزمایش، تنها حدود ۱۷٪ از افراد با جواب مثبت واقعاً بیمارند! علت، کوچک بودن احتمال پیشین ($1\%$) است: تعداد افراد سالمی که به اشتباه مثبت می‌شوند ($5\%$ از $99\%$) از تعداد بیماران واقعی بیشتر است.
\end{solution}
\end{example}

\begin{example}[حل دقیق مسئله مونتی هال با قانون بیز]
فرض کنید در ۱ را انتخاب کرده‌اید و مجری در ۳ را باز کرده است. احتمال اینکه جایزه پشت در ۲ باشد چقدر است؟ (فرض: اگر جایزه پشت در ۱ باشد، مجری با احتمال برابر یکی از درهای ۲ و ۳ را باز می‌کند.)
\begin{solution}
$C_i$ را رویداد «جایزه پشت در $i$» و $O_3$ را «مجری در ۳ را باز می‌کند» بگیرید. احتمال‌های پیشین $\Prob(C_i)=\frac13$ هستند. درست‌نمایی‌ها:
\[
\Prob(O_3\given C_1)=\tfrac12,\qquad \Prob(O_3\given C_2)=1,\qquad \Prob(O_3\given C_3)=0 .
\]
با احتمال کل: $\Prob(O_3)=\frac13(\frac12+1+0)=\frac12$. با قانون بیز:
\[
\Prob(C_2\given O_3)=\frac{1\times\frac13}{\frac12}=\frac23,\qquad \Prob(C_1\given O_3)=\frac{\frac12\times\frac13}{\frac12}=\frac13 .
\]
پس تغییر انتخاب احتمال برد را دو برابر می‌کند.
\end{solution}
\end{example}

\section{استقلال شرطی \lr{(Conditional Independence)}}

در موارد متعددی پیش می‌آید که دو رویداد $A$ و $B$ از یکدیگر مستقل نباشند:
\[
\Prob(A,B)\neq\Prob(A)\Prob(B),
\]
ولی در صورت دانستن رویداد سومی مانند $C$، به صورت شرطی از یکدیگر مستقل شوند.

\begin{definition}[استقلال شرطی]
دو رویداد $A$ و $B$ را به شرط $C$ مستقل گوییم (و می‌نویسیم $A\perp B\given C$) اگر
\[
\Prob(A,B\given C)=\Prob(A\given C)\,\Prob(B\given C).
\]
معادلاً (وقتی احتمال‌ها مثبت‌اند): $\Prob(A\given B,C)=\Prob(A\given C)$؛ یعنی وقتی $C$ را می‌دانیم، دانستن $B$ اطلاع بیشتری درباره $A$ نمی‌دهد.
\end{definition}

این اتفاق زمانی رخ می‌دهد که مسیر جریان اطلاعات (\lr{information flow}) از سمت رویداد $A$ به سمت $B$ (و یا بالعکس)، به واسطه دادن/گرفتن اطلاعاتی در مورد رویداد $C$ باشد:

\begin{center}
\begin{tikzpicture}[node distance=3cm,every node/.style={circle,draw,very thick,minimum size=1.1cm,font=\Large}]
\node (A) {$A$};
\node[right of=A] (C) {$C$};
\node[right of=C] (B) {$B$};
\draw[-{Stealth},very thick] (A) to[bend left=25] (C);
\draw[-{Stealth},very thick] (C) to[bend left=25] (A);
\draw[-{Stealth},very thick] (C) to[bend left=25] (B);
\draw[-{Stealth},very thick] (B) to[bend left=25] (C);
\end{tikzpicture}
\end{center}

اگر $C$ معلوم باشد، این «کانال» مسدود می‌شود و $A$ و $B$ دیگر از طریق آن اطلاعاتی رد و بدل نمی‌کنند.

\begin{warning}
\warntitle استقلال و استقلال شرطی هیچ‌کدام دیگری را نتیجه نمی‌دهند. ممکن است $A\perp B$ ولی $A\not\perp B\given C$ (مثلاً دو سکه مستقل و $C$=«دو سکه یکسان‌اند»؛ با دانستن $C$، نتیجه سکه اول، سکه دوم را کاملاً مشخص می‌کند). و ممکن است $A\not\perp B$ ولی $A\perp B\given C$ (مثال بعد).
\end{warning}

\begin{example}[دو سکه نامتقارن ــ صورت مسئله]\label{ex:twocoins}
فرض کنید در جعبه‌ای دو سکه تصادفی و نامتقارن گذارده شده‌اند که احتمال شیر آمدن در آنان به ترتیب ۰٫۲۵ و ۰٫۸ است. بر این اساس، آزمون تصادفی زیر را انجام می‌دهیم:
\begin{enumerate}
\item به صورت تصادفی و با احتمال‌های برابر، یکی از دو سکه را از جعبه انتخاب می‌کنیم.
\item این سکه را ۵ بار به صورت مستقل پرتاب کرده و نتایج را در گزارش $A$ ثبت می‌کنیم.
\item سپس، سکه را ۵ بار دیگر (به صورت مستقل از یکدیگر و بند ۲) پرتاب کرده و نتایج را در گزارش $B$ ثبت می‌کنیم.
\end{enumerate}
\textbf{سؤال:} در صورت دانستن نتیجه گزارش $A$، پیش‌بینی شما در مورد گزارش $B$ تغییر می‌کند یا خیر؟ چرا؟ فرض کنید که نتیجه گزارش $A$ برابر با $HHTHT$ شود. آنگاه احتمال اینکه در گزارش $B$ رشته $TTHTT$ را ببینیم چقدر است؟
\end{example}

پاسخ کامل این مثال را در ابتدای جلسه بعد خواهیم دید. پاسخ کیفی: \textbf{بله، تغییر می‌کند}؛ زیرا گزارش $A$ اطلاعاتی درباره اینکه کدام سکه انتخاب شده است به ما می‌دهد (سه شیر از پنج پرتاب، احتمال سکه دوم را بالا می‌برد) و این اطلاع، پیش‌بینی ما از $B$ را تغییر می‌دهد. با این حال، \emph{به شرط دانستن سکه}، $A$ و $B$ مستقل‌اند.

\section{تمرین‌های تکمیلی}

\begin{exercise}
دو تاس سالم پرتاب می‌کنیم. به شرط اینکه مجموع دو تاس ۸ باشد، احتمال اینکه دست‌کم یکی از تاس‌ها ۶ باشد چقدر است؟
\end{exercise}
\begin{solution}
حالت‌های با مجموع ۸: $(2,6),(3,5),(4,4),(5,3),(6,2)$ ــ پنج حالت هم‌احتمال. از این‌ها دو حالت شامل ۶ هستند. پس $\Prob=\frac25$. (در مقایسه، احتمال غیرشرطی وجود دست‌کم یک ۶ برابر $\frac{11}{36}\approx0.31$ است.)
\end{solution}

\begin{exercise}[دو فرزند]
خانواده‌ای دو فرزند دارد (هر فرزند مستقلاً با احتمال $\frac12$ پسر است). (الف) به شرط اینکه فرزند بزرگ‌تر پسر باشد، احتمال اینکه هر دو پسر باشند چقدر است؟ (ب) به شرط اینکه دست‌کم یکی پسر باشد چطور؟
\end{exercise}
\begin{solution}
فضای نمونه $\{BB,BG,GB,GG\}$ هم‌احتمال است (حرف اول: فرزند بزرگ‌تر).
(الف) شرط: $\{BB,BG\}$؛ پس $\Prob(BB\given\cdot)=\frac12$.
(ب) شرط: $\{BB,BG,GB\}$؛ پس $\Prob(BB\given\cdot)=\frac13$. تفاوت دو جواب نشان می‌دهد که «چه اطلاعاتی» دقیقاً داده شده، اهمیت دارد.
\end{solution}

\begin{exercise}[قانون بیز با سه خط تولید]
در مثال کارخانه، اگر محصولی معیوب باشد، احتمال اینکه از خط سوم آمده باشد چقدر است؟
\end{exercise}
\begin{solution}
با قانون بیز:
$\Prob(L_3\given D)=\frac{0.05\times0.2}{0.021}=\frac{0.010}{0.021}\approx0.476$.
به طور مشابه $\Prob(L_1\given D)=\frac{0.005}{0.021}\approx0.238$ و $\Prob(L_2\given D)=\frac{0.006}{0.021}\approx0.286$ (جمع سه احتمال پسین برابر ۱ است). خط سوم با اینکه فقط ۲۰٪ تولید را دارد، تقریباً نیمی از محصولات معیوب را تولید می‌کند.
\end{solution}

\begin{exercise}[آزمایش تکراری]
در مثال آزمایش پزشکی، اگر فردی دو بار (به صورت مستقل به شرط وضعیت بیماری) آزمایش دهد و هر دو بار مثبت شود، احتمال بیماری او چقدر است؟
\end{exercise}
\begin{solution}
می‌توان احتمال پسین مرحله اول ($0.167$) را به عنوان پیشین مرحله دوم به کار برد:
\[
\Prob(D\given ++)=\frac{0.99\times0.1667}{0.99\times0.1667+0.05\times0.8333}=\frac{0.165}{0.165+0.0417}\approx0.798 .
\]
همین جواب مستقیماً نیز به دست می‌آید: $\frac{0.99^2\times0.01}{0.99^2\times0.01+0.05^2\times0.99}\approx0.798$. (این‌جا از استقلال شرطی دو آزمایش به شرط وضعیت بیماری استفاده کردیم.)
\end{solution}
